\documentclass[11pt,reqno]{amsart}
\usepackage[T1]{fontenc}
\usepackage{lmodern}
\usepackage{microtype}
\usepackage{amsmath,amssymb,mathtools}
\usepackage{booktabs}
\usepackage{longtable}
\usepackage[margin=1.15in]{geometry}
\usepackage[colorlinks=true,linkcolor=blue,citecolor=blue,urlcolor=blue]{hyperref}
\newtheorem{theorem}{Theorem}[section]
\newtheorem{lemma}[theorem]{Lemma}
\newtheorem{proposition}[theorem]{Proposition}
\newtheorem{corollary}[theorem]{Corollary}
\newtheorem{conjecture}[theorem]{Conjecture}
\newtheorem*{mainthm}{Main Theorem}
\theoremstyle{definition}
\newtheorem{definition}[theorem]{Definition}
\newtheorem{remark}[theorem]{Remark}
\newtheorem{computation}[theorem]{Computation}
\newtheorem{assessment}[theorem]{Assessment}
\newtheorem{problem}{Problem}
\newcommand{\Z}{\mathbb Z}
\newcommand{\Q}{\mathbb Q}
\newcommand{\N}{\mathbb N}
\newcommand{\R}{\mathbb R}
\newcommand{\Zh}{\widehat{\Z}}
\newcommand{\Cl}{\operatorname{Cl}}
\newcommand{\ord}{\operatorname{ord}}
\newcommand{\lc}{\operatorname{lc}}
\newcommand{\leg}[2]{\left(\frac{#1}{#2}\right)}
\newcommand{\status}[1]{\textup{\textsc{#1}}}
\title[Erd\H{o}s--Straus for exponential families]{The Erd\H{o}s--Straus equation on exponential families beyond every Elsholtz--Tao covering}
\author{Anonymous}
\date{Internal draft (task O277), 19 October 2026; internal checks only, not externally refereed}
\emergencystretch=3em
\begin{document}
\begin{abstract}
For an exponential family $n_s=c_0\,2^s+k$ the residues of $n_s$ modulo any
fixed integer are eventually periodic in $s$, so one may try to prove the
Erd\H{o}s--Straus conjecture $4/n=1/x+1/y+1/z$ for all members at once by a
finite covering of the exponents with Elsholtz--Tao identity classes and
divisor classes, in the manner of Sierpi\'nski's coverings for $k\cdot2^n+1$.
When $c_0+k=4^J$ and $k\equiv1\pmod 8$ this is impossible: along
$s\equiv0\pmod L$ the members imitate the even square $4^J$ modulo every
odd number, and an odd-square (Schinzel--Yamamoto) argument shows that no
finite set of such classes covers them.  We show that these
obstructed families can nevertheless be settled completely, by adding a few
identity classes whose coordinates are polynomials in a root $Y=2^{s/D}$ of
the exponential rather than in $n$.  Their divisors are forced by an
algebraic factorisation along the family (Sophie Germain's identity
$4X^4+1=(2X^2+2X+1)(2X^2-2X+1)$, cyclotomic factors of $Y^h+w^h$, or a
binomial $n+4a^2d$), in the spirit of Izotov's construction for Sierpi\'nski
numbers.  We prove the Erd\H{o}s--Straus conjecture for every
$n=3\cdot2^s+1$, every $63\cdot2^s+1$ and every $7\cdot4^s+9$ (all three
provably beyond every finite Elsholtz--Tao covering), and for every
$4^s+k$ with $k\in\{9,57,81,105,129,153,177,201,225,249\}$; for
$15\cdot2^s+1$ we prove it for every $s\not\equiv0\pmod{60}$.  All proofs
are explicit identity certificates; most were found by computer search and
are re-verified by a short self-contained script.  On the negative side we
prove that no identity class of constant gap $4x-n$ or constant ratio $x/n$,
polynomial in any root $2^{s/D}$, reaches the unit-square point of
$15\cdot2^s+1$, $31\cdot4^s+33$, $7\cdot4^s+57$, $63\cdot4^s+1$,
$47\cdot4^s+17$ or $55\cdot4^s+9$, and we close one branch of the
moving-gap binomial menu at $15\cdot2^s+1$.  The families are thin; nothing
here bears on the Erd\H{o}s--Straus conjecture for primes in general.
\end{abstract}
\maketitle
\tableofcontents

\section{Introduction}\label{sec:intro}

The Erd\H{o}s--Straus conjecture (ES) asserts that for every integer $n\ge2$
\begin{equation}\label{eq:ES}
 \frac4n=\frac1x+\frac1y+\frac1z
\end{equation}
has a solution in positive integers.  If $n=ab$ and \eqref{eq:ES} is
solvable for $a$, it is solvable for $n$ (multiply $x,y,z$ by $b$), so it
suffices to treat primes.  The conjecture has been verified numerically
up to $10^{17}$ by Salez \cite{Salez} and up to $10^{18}$ by Mihnea and
Dumitru \cite{MD}.

\subsection*{Identity classes and their limits}
The classical tool is the polynomial identity on a residue class.  Mordell
\cite{Mordell} showed that every $n$ outside the six classes
$1,121,169,289,361,529\pmod{840}$ (the squares modulo $840$) is covered by
such identities, and Elsholtz and Tao \cite[Prop.~1.9]{ET} essentially
classified the residue classes that are \emph{solvable by polynomials}: up
to finitely many primes they are unions of classes of seven explicit
families (Table~\ref{tab:ET}).  A primitive class whose residue is a perfect
square is never solvable by polynomials (Mordell, Schinzel
\cite{Mordell,Schinzel}; it also follows from the vanishing of Type~I and
Type~II solutions at odd squares, \cite[Prop.~1.6]{ET}).  Elsholtz and Tao
remark that this ``rules out \dots\ a finite set of covering congruence
strategies'' for ES \cite[p.~6]{ET}.  A companion note \cite{CovNote}
studies this obstruction for sets of primes from the profinite point of
view.

\subsection*{Exponential families}
Here we look at thin families
\[
  \mathcal F(c_0,b,k)=\{\,n_s=c_0\,b^s+k : s\ge0\,\},\qquad b\in\{2,4\},
\]
with $c_0,k$ positive integers.  The residue of $n_s$ modulo a fixed $M$ is
eventually periodic in $s$, so a finite list of identity classes may cover
\emph{all} exponents, exactly as Sierpi\'nski \cite{Sierpinski} proved that
$k\cdot2^n+1$ is composite for all $n$ by covering the exponents $n$ with
finitely many congruences (see \cite[\S1]{FFK}).  We call such a list an
\emph{exponent covering}.  For many families it exists and is small
(Theorem~\ref{thm:nine}); the companion work \cite{ExpCov} certified,
for instance, ES for every $315\cdot2^m+1$ and every $15015\cdot2^m+1$.

The obstruction is a \emph{unit-square point}.  If $c_0$ is odd,
$k\equiv1\pmod8$ and $c_0+k=4^J$, then for $s\equiv0\pmod L$ with $L$
highly divisible, $n_s\equiv c_0+k=4^J$ modulo every odd number of the form
$2^L-1$ and $n_s\equiv k\pmod{2^s}$; so $n_s$ is congruent to an odd
perfect square modulo every fixed modulus, and the odd-square argument
shows that no finite set of Elsholtz--Tao classes and divisor classes
covers the family (Theorem~\ref{thm:sterile}).  Examples are
$3\cdot2^s+1$ (the Proth numbers with $k=3$), $15\cdot2^s+1$,
$63\cdot2^s+1$, $7\cdot4^s+9$, $31\cdot4^s+33$.

\subsection*{Results}
The point of this note is that the obstruction concerns identities that
are polynomial in $n$, not ES itself.  Along a residue class of exponents
the integer $n_s$ is a polynomial in a root $Y=2^{(s-s_0)/D}$, and an
identity polynomial in $Y$ is again valid on a whole residue class of $s$.
Such an identity can use divisors that move with $n$ and are forced by an
algebraic factorisation of $Y$-polynomials, and these can reach the
unit-square point.

\begin{mainthm}
\begin{enumerate}
\item[(a)] \textup{(Theorems~\ref{thm:proth3}, \ref{thm:63}, \ref{thm:c7})}
The Erd\H{o}s--Straus conjecture holds for every $n=3\cdot2^s+1$, every
$n=63\cdot2^s+1$ and every $n=7\cdot4^s+9$ ($s\ge0$).  None of these
families admits a finite covering by Elsholtz--Tao classes and divisor
classes (Theorem~\ref{thm:sterile}).
\item[(b)] \textup{(Theorems~\ref{thm:105}, \ref{thm:nine})} It holds for
every $n=4^s+k$, $s\ge0$, for
$k\in\{9,57,81,105,129,153,177,201,225,249\}$.
\item[(c)] \textup{(Proposition~\ref{prop:15half})} It holds for every
$n=15\cdot2^s+1$ with $60\nmid s$.
\item[(d)] \textup{(Theorem~\ref{thm:W}, Corollary~\ref{cor:ratio})} For
$(c_0,k)\in\{(15,1),(31,33),(7,57),(63,1),(47,17),(55,9)\}$ no identity
family that is polynomial in a root $Y=2^{s/D}$ (any $D$) and has constant
gap $4x-n$ or constant ratio $x/n$ is valid on a neighbourhood
$\{s\equiv0\pmod L,\ s\ge s_0\}$ of the unit-square point.  Hence
these families have no finite covering by Elsholtz--Tao classes, divisor
classes and such families.
\end{enumerate}
\end{mainthm}

Parts (a)--(c) are \status{proved}; the certificates of (a) for
$63\cdot2^s+1$ and $7\cdot4^s+9$, of (b) and of (c) were found by computer
search and the proofs are \emph{computer-assisted} in the sense that a
finite coverage statement (a list of at most $9180$ residues) is checked by
machine; every identity is printed in the text and can be checked by hand.
The proof for $3\cdot2^s+1$ (Theorem~\ref{thm:proth3}) is short and fully
hand-checkable.  Part (d) is \status{proved} up to a finite enumeration
(Computation~\ref{comp:W}).  In \S\ref{sec:fifteen} we also prove that one of
the two branches of the moving-gap menu that produced (a) for
$63\cdot2^s+1$ is empty for $15\cdot2^s+1$ (Theorem~\ref{thm:iiad3}); the
other branch, and ES for $15\cdot2^s+1$ with $60\mid s$, remain open.

\subsection*{The mechanism, and a precedent}
In all three families of (a) the unit-square point is reached by a single
kind of move.  For $3\cdot2^s+1$ with $4\mid s-4$ put $X=2^{(s-4)/4}$, so
$n=48X^4+1$ and $x:=(n+11)/4=3(4X^4+1)$.  Sophie Germain's identity
\[
  4X^4+1=(2X^2+2X+1)(2X^2-2X+1)
\]
supplies divisors of $x$ that are not powers of fixed primes, and the
elementary window lemma (Lemma~\ref{lem:window}) turns a congruence
between them modulo the \emph{window} $m=4x-n=11$ into a solution of
\eqref{eq:ES}.  The congruence holds on whole residue classes of $s$,
including a neighbourhood of the unit-square point.  This is the ES
counterpart of Izotov's observation \cite{Izotov} (see \cite[pp.~3--4]{FFK})
that a Sierpi\'nski-type covering may omit one congruence class of
exponents if the family factors algebraically there, by the same
factorisation $4x^4+1=(2x^2+2x+1)(2x^2-2x+1)$.  For $4^s+105$ the closing
class uses $u^3+1=(u+1)(u^2-u+1)$, $u=4^{(s-4)/3}$; for $63\cdot2^s+1$ it is
a Type~II family with a \emph{moving} gap $4x-n=3(Y^2-3Y+9)$ built from
$Y^3+27=(Y+3)(Y^2-3Y+9)$ (\S\ref{sec:forced}).  Part (d) shows that the
constant-gap version of the mechanism has a precise 2-adic limit: it can
reach a unit-square point only through a Sophie Germain piece whose value
at the limit has odd 2-adic valuation (Lemma~\ref{lem:limitjacobi}), and an
explicit finite check shows that no such piece exists for the six listed
families.

\subsection*{Scope}
Each family has $O(\log X)$ members up to $X$, so these theorems say
nothing about ES for primes in general; a finite set of thin identity
classes misses almost all primes of any residue class (see
\cite[Prop.~5.1]{MenuSterility}).  The members $s\le 50$ or so are already
covered by the numerical verifications \cite{Salez,MD}.  Each family lies
in a finite union of residue classes, and several of our families are
settled by Mordell's identities for most $s$; we say where the new content
lies.  We found no prior work on ES for exponential families (searches of
zbMATH Open for all documents on the Erd\H{o}s--Straus equation, and web
searches; see \cite{Audit}); this is not a claim of priority.

\subsection*{Labels}
\status{Proved} means that a complete proof is given here.
\status{Computer-assisted} means that a finite verification (always of a
statement about finitely many explicit residues and identities) is done by
the script \texttt{scripts/c1\_expfam\_verify.py}
(Appendix~\ref{app:script}), which uses only exact arithmetic and imports
nothing from the search code; the original certificates were in addition
checked by two independent reviewers' programs each.  \status{Evidence}
means a computation that supports but does not prove a statement.
\status{Open} and \status{Conjecture} are not results.

\subsection*{Provenance}
This note consolidates internally reviewed work of a research notebook:
\cite{ExpCov} (exponent coverings, the unit-square criterion),
\cite{SignFields} ($4^s+k$), \cite{Sparks6} ($4^s+105$),
\cite{MenuSterility} ($3\cdot2^s+1$), \cite{Exp105} ($7\cdot4^s+9$ and a
second proof for $3\cdot2^m+1$), \cite{ExpOpen} ($63\cdot2^s+1$ and
Theorem~\ref{thm:W}) and \cite{Exp15} ($15\cdot2^s+1$).  Several proofs are
re-derived and expanded here; Proposition~\ref{prop:15half} is new in this
note.

\section{Exponent coverings and unit-square points}\label{sec:criterion}

\subsection{The Elsholtz--Tao classes}
We use the seven families of \cite[Prop.~1.9]{ET} in the notation of
\cite[\S2]{CovNote}.  Each is indexed by a triple $\kappa$ of positive
integers subject to a side condition and defines a set of residue classes
modulo $M(\kappa)$ (Table~\ref{tab:ET}).  On its class each family carries
an explicit solution
\[
 \pi^{I}(a,b,c,d)=(abdn,\ acd,\ bcd)\quad\text{(I1--I4)},\qquad
 \pi^{II}(a,b,c,d)=(abd,\ acdn,\ bcdn)\quad\text{(II1--II3)},
\]
where three of $a,b,c,d$ are given by $\kappa$ or by $\kappa$ alone, and the
remaining ones are the polynomials in $n$ listed in the table.
\begin{table}[ht]
\centering\small
\begin{tabular}{lllll}
\toprule
family & $\kappa$ & side condition & class & remaining coordinates\\
\midrule
I1 & $a,d,f$ & $f\mid 4a^2d+1$ & $n\equiv-f \pmod{4ad}$ & $e=\frac{4a^2d+1}{f}$, $b=\frac{ne+1}{4ad}$, $c=\frac{n+f}{4ad}$\\
I2 & $a,c,f$ & $(4ac,f)=1$ & $n\equiv-f\ (4ac)$, $n\equiv -c/a\ (f)$ & $d=\frac{n+f}{4ac}$, $b=\frac{na+c}{f}$\\
I3 & $c,d,f$ & $(4cd,f)=1$ & $n\equiv-f\ (4cd)$, $f\mid n^2+4c^2d$ & $a=\frac{n+f}{4cd}$, $b=\frac{n^2+4c^2d+nf}{4cdf}$\\
I4 & $a,b,e$ & $e\mid a+b$, $(e,4ab)=1$ & $n\equiv-1/e \pmod{4ab}$ & $c=\frac{a+b}{e}$, $d=\frac{ne+1}{4ab}$\\
II1 & $a,b,e$ & $e\mid a+b$, $(e,4ab)=1$ & $n\equiv-e\pmod{4ab}$ & $c=\frac{a+b}{e}$, $d=\frac{n+e}{4ab}$\\
II2 & $a,d,f$ & $4ad\mid f+1$ & $n\equiv-4a^2d\pmod f$ & $c=\frac{f+1}{4ad}$, $b=\frac{nc+a}{f}$\\
II3 & $a,d,e$ & $(4ad,e)=1$ & $n\equiv-4a^2d-e\pmod{4ade}$ & $c=\frac{n+4a^2d+e}{4ade}$, $b=\frac{n+e}{4ad}$\\
\bottomrule
\end{tabular}
\caption{The Elsholtz--Tao families \cite[Prop.~1.9]{ET}.}\label{tab:ET}
\end{table}

\begin{lemma}[\status{proved}; {\cite[\S\S2,10]{ET}, \cite[Lemma~2.1]{CovNote}}]\label{lem:solve}
If $\kappa$ is admissible and the integer $n\ge1$ lies in its class, the
coordinates of Table~\ref{tab:ET} are positive integers and $\pi^{I}$
\textup{(}resp.\ $\pi^{II}$\textup{)} is a solution of \eqref{eq:ES}.
Moreover $4/n=1/(abdn)+1/(acd)+1/(bcd)$ is equivalent to
$4abcd=n(a+b)+c$, and $4/n=1/(abd)+1/(acdn)+1/(bcdn)$ to $4abcd=nc+a+b$.
\end{lemma}

The last two equivalences are immediate on clearing denominators.  We call
the class of $\kappa$ an \emph{ET class}.  A \emph{divisor class} is
$D_q=\{n\equiv0\pmod q\}$ for an odd prime $q$; on it
$(x,y,z)=\frac nq(x_q,y_q,z_q)$ for a fixed solution $(x_q,y_q,z_q)$ of
$4/q=1/x+1/y+1/z$.  (All divisor classes used below have $q\le 541$, and
the script records explicit solutions.)

\subsection{Odd squares}
The following is the classical quadratic-reciprocity obstruction of
Schinzel and Yamamoto \cite{Yamamoto}, in the form proved by Elsholtz and Tao
for Type~I \cite[\S4]{ET} (they omit the Type~II case).  We give both
cases, since every negative statement below rests on it.

\begin{lemma}[\status{proved}]\label{lem:T1}
Let $n$ be an odd perfect square.
\begin{enumerate}
\item[(i)] There are no positive integers $a,b,c,d$ with
$4abcd=n(a+b)+c$ and $\gcd(c,n)=1$.
\item[(ii)] There are no positive integers $a,b,c,d$ with
$4abcd=nc+a+b$ and $\gcd(abd,n)=1$.
\end{enumerate}
\end{lemma}

\begin{proof}
Note $n\equiv1\pmod8$, and write $q$ for the odd part of $ab$.  We use the
Jacobi symbol and the reciprocity law for odd positive coprime $q,e$:
$\leg qe=\leg eq(-1)^{\frac{q-1}2\frac{e-1}2}$.

(i) Since $c\mid n(a+b)$ and $\gcd(c,n)=1$, $e:=(a+b)/c$ is a positive
integer, and $4abd=ne+1$.  Hence $\gcd(e,ab)=1$, $\gcd(n,q)=1$, and
$ne\equiv-1\pmod 4$, so $e\equiv3\pmod4$.  Modulo $q$, $e\equiv-n^{-1}$, and
$n$ is a square, so $\leg eq=\leg{-1}q$.  By reciprocity with $e\equiv3\pmod
4$, $\leg qe=\leg eq\leg{-1}q=1$.  On the other hand $b\equiv-a\pmod e$ and
$\gcd(a,e)=1$, so $\leg{ab}e=\leg{-a^2}e=\leg{-1}e=-1$.  Writing
$ab=2^vq$ we get $\leg 2e^v=-1$, so $v\ge1$.  Then $8\mid4abd=ne+1$, so
$e\equiv7\pmod8$ and $\leg2e=1$, a contradiction.

(ii) Put $e:=4abd-n$.  Then $ce=4abcd-nc=a+b>0$, so $e\ge1$.  Since
$\gcd(ab,n)=1$ we have $\gcd(n,q)=1$; $e\equiv-n\pmod q$ gives
$\leg eq=\leg{-1}q$; and $e\equiv-n\equiv3\pmod4$.  As in (i),
$\leg qe=1$.  A prime dividing $e$ and $a$ (or $e$ and $b$) divides
$4abd-e=n$, which is excluded, so $\gcd(ab,e)=1$, and $b\equiv-a\pmod e$
gives $\leg{ab}e=-1$.  Hence $ab$ is even, $8\mid 4abd$,
$e\equiv-n\equiv7\pmod8$ and $\leg2e=1$, a contradiction.
\end{proof}

\begin{corollary}[\status{proved}]\label{cor:T1class}
Let $\kappa$ be an admissible triple of one of the seven families, let
$P_\kappa$ be the product of those coordinates among $a,b,c,d,e,f$ that do
not depend on $n$ \textup{(}the entries of $\kappa$, together with $e$ for
I1, $c$ for I4, II1 and II2\textup{)}, and let $t$ be an odd integer with
$\gcd(t,P_\kappa)=1$.  Then $t^2$ does not lie in the class of $\kappa$.
\end{corollary}

\begin{proof}
Suppose it does and put $n=t^2$.  By Lemma~\ref{lem:solve} the coordinates
$a,b,c,d$ are positive integers with $4abcd=n(a+b)+c$ (Type~I) or
$4abcd=nc+a+b$ (Type~II).  For Type~I we need $\gcd(c,n)=1$: in I2, I3 and
I4 the coordinate $c$ divides $P_\kappa$, and in I1 the relation $4ad\,c=n+f$
gives $\gcd(c,n)\mid f$.  For Type~II we need $\gcd(abd,n)=1$: the
constant coordinates divide $P_\kappa$, and for the others, in II1
$4ab\,d=n+e$ gives $\gcd(d,n)\mid e$, in II2 $fb=nc+a$ gives
$\gcd(b,n)\mid a$, and in II3 $4ad\,b=n+e$ gives $\gcd(b,n)\mid e$.  In
every case the relevant gcd divides $P_\kappa$ and is therefore $1$, and
Lemma~\ref{lem:T1} gives a contradiction.
\end{proof}

\subsection{Rules and exponent coverings}
Fix a family $\mathcal F(c_0,b,k)$.  For an integer $M=2^aM'$ with $M'$ odd,
\begin{equation}\label{eq:period}
 n_{s+\ell}\equiv n_s \pmod{M}\qquad\text{whenever } \ord_{M'}(b)\mid \ell \text{ and } b^s\equiv0\pmod{2^a},
\end{equation}
since $n_{s+\ell}-n_s=c_0b^s(b^\ell-1)$.  So for every ET or divisor class
$C$ the set $\{s\ge s_1: n_s\in C\}$ is, for $s_1$ large, a union of
residue classes modulo $\ord_{M'}(b)$.  More generally:

\begin{definition}\label{def:rule}
A \emph{rule} for $\mathcal F(c_0,b,k)$ consists of integers $s_0\ge0$,
$\sigma\ge1$, $\omega\ge1$ such that $D:=\sigma\log_2 b/\omega$ is an integer,
and of $x,y,z\in\Q[Y]$ with
\[
 4xyz=n(Y)\,(xy+yz+zx),\qquad n(Y):=c_0b^{s_0}Y^D+k .
\]
For $s=s_0+\sigma t$ ($t\ge0$) put $Y=2^{\omega t}$; then $n(Y)=n_s$.  The
rule \emph{applies} at $s$ if $x,y,z$ take positive integer values there.
An \emph{exponent covering} of the family is a finite list of rules such
that every $s\ge0$ at which no rule applies is treated separately.
\end{definition}

ET and divisor classes are rules with $s_0=0$, $\sigma=1$, $Y=b^s$ and
$n(Y)=c_0Y+k$: their coordinates are polynomials in $n$.  The point of
this note is to use rules with $D\ge2$.

\begin{lemma}[\status{proved}]\label{lem:ruleperiod}
Let a rule be given, and write each of $x,y,z$ as $R_1(Y)/N$ with
$R_1\in\Z[Y]$ and $N=2^aN'$, $N'$ odd.  There are $t^*\ge0$ and $T\ge1$
such that for $t\ge t^*$ whether the rule applies at $s=s_0+\sigma t$
depends only on $t\bmod T$.  One may take $T$ the least common multiple of
the orders $\ord_{N'}(2^\omega)$ over the three coordinates, and $t^*$ such
that $\omega t^*\ge a$ for all three and all numerators are positive on
$[2^{\omega t^*},\infty)$.
\end{lemma}

\begin{proof}
For $\omega t\ge a$ we have $R_1(2^{\omega t})\equiv R_1(0)\pmod{2^a}$, and
$R_1(2^{\omega t})\bmod N'$ is periodic in $t$ with period
$\ord_{N'}(2^\omega)$.  Positivity holds for all $t\ge t^*$ by the choice of
$t^*$.
\end{proof}

So the set of exponents at which a rule applies is, up to finitely many
$s$, a finite union of residue classes $s\equiv u\pmod{\sigma T}$.  The
script of Appendix~\ref{app:script} computes $T$, $t^*$ and the exact set of
admissible residues $t\bmod T$ for every rule, by evaluating one full period.
In profinite terms, these sets are clopen subsets of the exponent space
$\Zh$, and a finite covering of all large $s$ by them covers all of $\Zh$,
because the integers $s\ge s_1$ are dense in $\Zh$.  We use this only in the
following elementary form.

\begin{lemma}[\status{proved}]\label{lem:compact}
If finitely many sets $E_i=\{s\equiv u_i\pmod{P_i},\ s\ge s_i\}$ cover all
sufficiently large $s$, then for some $i$ we have $u_i\equiv0\pmod{P_i}$, so
$E_i$ contains every sufficiently large multiple of $P_i$.
\end{lemma}

\begin{proof}
Let $P=\operatorname{lcm}P_i$.  A large multiple $s$ of $P$ lies in some
$E_i$, so $u_i\equiv s\equiv0\pmod{P_i}$.
\end{proof}

\subsection{Unit-square points}
\begin{theorem}[\status{proved}; {\cite[Thm.~1.4, Cor.~1.5]{ExpCov}, \cite[Lemma~3.1]{ExpNearSq}}]\label{thm:sterile}
Let $c_0$ be odd, $k\equiv1\pmod 8$ and $c_0+k=4^J$.  For every $L\ge1$ and
$s_1\ge0$ there is no finite set of ET classes and divisor classes
containing all $n_s=c_0\,2^s+k$ with $s\equiv0\pmod L$, $s\ge s_1$.
\end{theorem}

\begin{proof}
Suppose $C_1,\dots,C_r$ are such classes, of moduli $M_1,\dots,M_r$, and
let $M=\operatorname{lcm}M_i=2^aM'$ with $M'$ odd.  For
$s\in S:=\{s\ge\max(a,s_1): L\ord_{M'}(2)\mid s\}$ we have $n_s\equiv
c_0+k=4^J\pmod{M'}$ and $n_s\equiv k\pmod{2^a}$, so $\rho:=n_s\bmod M$ does
not depend on $s\in S$, and $\rho\in C_i$ for some $i$.  If $C_i=D_q$, then
$q\mid M'$ and $q\mid\rho\equiv4^J\pmod q$, impossible.  So $C_i$ is the
class of some $\kappa$.  Since $k\equiv1\pmod8$, $k$ is a square in
$\Z_2^\times$; choose by the Chinese remainder theorem an odd integer $t$
with $t\equiv2^J\pmod{M'}$, $t^2\equiv k\pmod{2^a}$, and $t\equiv1$ modulo
every prime of $P_\kappa$ not dividing $M$.  Then $t^2\equiv\rho\pmod M$, so
$t^2$ lies in the class of $\kappa$, and $\gcd(t,P_\kappa)=1$ (the primes of
$P_\kappa$ dividing $M'$ do not divide $t\equiv2^J$, and $t$ is odd).  This
contradicts Corollary~\ref{cor:T1class}.
\end{proof}

The same proof works for the subfamilies $s\equiv0\pmod 2$, i.e.\ for
$c_0\cdot4^{s'}+k$; so $3\cdot2^s+1$, $15\cdot2^s+1$, $63\cdot2^s+1$,
$7\cdot4^s+9$, $15\cdot4^s+1$, $23\cdot4^s+41$, $31\cdot4^s+33$,
$7\cdot4^s+57$, $63\cdot4^s+1$, $47\cdot4^s+17$, $55\cdot4^s+9$ admit no finite
ET/divisor covering.  The limit point behind the proof is $\mu=0$ in the
exponent space $\Zh$: there the $q$-adic limit of $n_s$ is $4^J$ for every
odd $q$ and the $2$-adic limit is $k$, a unit square at every place.  By
\cite[Prop.~1.9]{ET} the statement extends to all primitive classes solvable
by polynomials (cf.\ \cite[Prop.~3.1]{CovNote}); we do not need this.

The integer points of $\Zh$ that are unit-square points are classified as
follows.  Write $\nu(\ell)$ for the point of $\prod_q\Z_q$ with odd
components $c_0b^\ell+k\in\Z[1/2]$ and $2$-component $k$.

\begin{proposition}[\status{proved}; {\cite[Cor.~1.5]{ExpCov}, \cite[Lemma~3.1]{ExpNearSq}, \cite[Lemma~1.2]{SignFields}}]\label{prop:intpoints}
Let $k\equiv1\pmod 8$, $b\in\{2,4\}$, and $c_0$ odd.  For $\ell\in\Z$,
$\nu(\ell)$ is a square of a unit at every prime iff $\ell=0$ and
$c_0+k=4^J$ for some $J\ge1$.  For $c_0\equiv2\pmod 4$ and $b=4$ there is no
such $\ell$.
\end{proposition}

\begin{proof}
Let $\ell\ge0$ and $N=c_0b^\ell+k$.  A positive integer that is a square in
$\Z_q$ for all odd $q$ is a perfect square (an integer that is not a
square is a non-residue modulo infinitely many primes), and it is a unit at
all odd $q$ iff it is a power of $2$; so $N=4^J$.  For $\ell\ge1$ and $c_0$
odd $N$ is odd and $>1$; so $\ell=0$, $c_0+k=4^J$.  If
$\ell=-a<0$, the odd components are $(c_0+kb^{a})/b^{a}$; this is a unit
square everywhere iff $b^a(c_0+kb^a)$ is a perfect square with
$c_0+kb^a$ a power of $2$; but $c_0+kb^a$ is odd and $>1$.  For
$c_0\equiv2\pmod4$, $b=4$: $N=c_0b^\ell+k$ is odd for $\ell\ge1$, and
$c_0+k\equiv3\pmod 4$ is not a power of $4$; the case $\ell<0$ is as
before ($c_0+k4^a$ is odd).
\end{proof}

\begin{remark}[non-integer points; \status{open}]\label{rem:EH}
Whether the closure of a family can meet the unit squares at a
non-integer point of $\Zh$ is an exponential local--global question in the
style of Skolem's conjecture and the Hasse-type principle of
Bert\'ok and Hajdu \cite{BH}; we call it (EH) as in \cite[\S3]{ExpCov}.  It
is open in general.  For $4^s+105$ a finite screen shows that the closure
avoids the unit squares (\cite[\S1.4]{SignFields}: for every residue
$s\bmod360$ some prime $q<3\cdot10^6$ with $\ord_q(4)\mid360$ divides
$n_s$ or has $\leg{n_s}q=-1$; a certified finite computation).  Beyond the
unit squares, the sterile set of the ET menu can contain non-square points;
none is known on any of the families studied here.
\end{remark}

\section{Windows with a forced budget}\label{sec:forced}

\subsection{The window lemma}
The following is the elementary criterion for $m/N=1/y+1/z$ (Rav
\cite{Rav}; we include the proof), specialised to the first step
$4/n=1/x+m/(nx)$.

\begin{lemma}[window lemma; \status{proved}]\label{lem:window}
Let $n,x$ be positive integers with $m:=4x-n>0$, and put $N:=nx$.
\begin{enumerate}
\item[(a)] If $d_1,d_2$ are positive divisors of $N$ with
$m\mid d_1+d_2$, then
\[
 y=\frac{N(d_1+d_2)}{m\,d_2},\qquad z=\frac{N(d_1+d_2)}{m\,d_1}
\]
are positive integers and $4/n=1/x+1/y+1/z$.  Conversely every solution
with this $x$ arises in this way.
\item[(b)] If $\gcd(m,x)=1$, $d\mid x^2$ and $m\mid x+d$, then
$(x,\ n(x+d)/m,\ nx(x+d)/(md))$ is a solution in positive integers.
\end{enumerate}
\end{lemma}

\begin{proof}
(a) $N/d_2$ and $(d_1+d_2)/m$ are integers, so $y\in\N$, and likewise $z$;
and $1/y+1/z=\frac{m(d_1+d_2)}{N(d_1+d_2)}=\frac mN=\frac4n-\frac1x$.
Conversely, if $m/N=1/y+1/z$, let $g=\gcd(y,z)$, $y=gy'$, $z=gz'$; then
$mgy'z'=N(y'+z')$ and $\gcd(y'z',y'+z')=1$, so $y'z'\mid N$; with
$d_1:=y'\cdot N/(y'z')=N/z'$ and $d_2:=N/y'$ we get $d_1+d_2=N(y'+z')/(y'z')=mg$,
and $y=N(d_1+d_2)/(md_2)$, $z=N(d_1+d_2)/(md_1)$.
(b) $y\in\N$ since $m\mid x+d$.  For $z$ we show $d\mid x(x+d)/m$.  Let $p$
be a prime, $\beta=v_p(x)$, $\gamma=v_p(d)\le2\beta$.  If $\gamma>0$ then
$\beta>0$ and $p\nmid m$, so $v_p(x(x+d)/m)=\beta+v_p(x+d)$, which is
$\ge\beta+\gamma\ge\gamma$ if $\gamma<\beta$, $\ge2\beta\ge\gamma$ if
$\gamma=\beta$, and $=2\beta\ge\gamma$ if $\gamma>\beta$.  The identity is
$\frac1y+\frac1z=\frac{m}{n(x+d)}\bigl(1+\frac dx\bigr)=\frac m{nx}$.
\end{proof}

We call $m$ the \emph{window} or \emph{gap} and $x+d$, resp.\ the pair
$(d_1,d_2)$, the \emph{budget}.  For a single integer $n$ the lemma is a
lottery: one needs divisors of $nx$ in a prescribed class modulo $m$.  On an
exponential family it becomes a theorem on whole residue classes of $s$ as
soon as the divisors are \emph{forced} by an algebraic identity.

\subsection{Forced budgets}
Let $\mathcal F(c_0,b,k)$ be a family and $m\equiv-k\pmod 4$ a positive
integer.  On a residue class $s=s_0+\sigma t$ the integer $x=(n_s+m)/4$ is
a polynomial in $Y=2^{\omega t}$ with rational coefficients, and so is
$N=nx$.  Every factorisation of $N$ in $\Q[Y]$ gives candidate divisors
$d_1(Y),d_2(Y)$; if $m\mid d_1(Y)+d_2(Y)$ for one $t$ in a residue class
modulo the period of Lemma~\ref{lem:ruleperiod}, it holds for the whole
class, and Lemma~\ref{lem:window} produces a rule.  The algebraic
factorisations used below are:
\begin{enumerate}
\item[(SG)] \emph{Sophie Germain}: $4Y^4+w^4=(2Y^2+2Yw+w^2)(2Y^2-2Yw+w^2)$
and $Y^4+4w^4=(Y^2+2Yw+2w^2)(Y^2-2Yw+2w^2)$.  For $x=c_04^{j-1}(4^r+w^4)$
with $r=s-j$ odd one puts $Y=2^{(r-1)/2}$, $4^r=4Y^4$.
\item[(CY)] \emph{cyclotomic}: for odd $h$,
$Y^h+w^h=\prod_{e\mid h}w^{\varphi(2e)}\Phi_{2e}(Y/w)$; for
$x=c_04^{j-1}(4^r+w^h)$ with $h\mid r$ one puts $Y=4^{r/h}$.
\item[(BI)] \emph{binomial, moving gap}: for constants $a,d$, the Type~II
identity $4abcd=cn+a+b$ is equivalent to
\begin{equation}\label{eq:binom}
 (4abd-n)(4acd-1)=n+4a^2d ,
\end{equation}
because the left side equals $4ad(4abcd-b-cn)+n$.  So every factorisation
$n+4a^2d=f\,f'$ in $\Q[Y]$ with $b:=(f+n)/(4ad)$, $c:=(f'+1)/(4ad)$
integer-valued and positive on a class gives a rule
$(x,y,z)=(abd,acdn,bcdn)$.  Its gap $4x-n=f$ moves with $s$.
\end{enumerate}
In (SG) and (CY) the gap $m$ is constant and the budget moves; in (BI) the
gap itself is an algebraic factor.  The pieces available in (SG), (CY) are
the odd primes of $c_0$, the polynomial $n(Y)$ itself, and the algebraic
factors of $x$; powers of $2$ up to $2^{v_2(c_04^{j-1})}$ may be added.

\begin{remark}[no conflict with the odd-square obstruction]\label{rem:T1window}
A rule with $D\ge2$ applies on a set of exponents $\{s\equiv u\pmod P\}$.
This set of integers $n_s$ is contained in a residue class of $n$ (for
instance modulo $2^P-1$), which by Theorem~\ref{thm:sterile} may well be
square-like.  There is no contradiction: Lemma~\ref{lem:T1} concerns
identities whose coordinates are polynomials in $n$, while a rule is
polynomial in $Y=n_s$'s root, $Y=((n-k)/(c_0b^{s_0}))^{1/D}$, and is only
asserted on the thin family.  If a member of the class is a square, the
rule simply solves it.
\end{remark}

\subsection{Specialisation at a unit-square point}
Let $c_0+k=4^J$ and consider rules through the neighbourhood
$\{s\equiv0\pmod L\}$ of $\mu=0$ for $n_s=c_0\,2^s+k$.  After the
substitution $Y=2^{s/D}$, such $s$ have $Y\equiv1$ modulo any fixed odd
number and $Y\equiv0$ modulo any fixed power of $2$ once $L$ and $s$ are
large.  This gives a convenient criterion.

\begin{lemma}[specialisation; \status{proved}; {\cite[Lemma~2.1]{MenuLadder}}]\label{lem:special}
Let $R\in\Q[Y]$ and $D\ge1$.  Then $R(2^{s/D})\in\Z$ for all $s$ in some set
$\{s\equiv0\pmod L,\ s\ge s_1\}$ \textup{(}$D\mid L$\textup{)} if and only if
$R(1)\in\Z[1/2]$ and $R(0)\in\Z_{(2)}$.  More precisely, for every prime
power $p^e$: if $R(2^{s/D})\in p^e\Z_{(p)}$ for all $s$ in such a set, then
$R(1)\in p^e\Z_p$ when $p$ is odd and $R(0)\in 2^e\Z_2$ when $p=2$.
\end{lemma}

\begin{proof}
Write $R=R_1/N$ with $R_1\in\Z[Y]$.  For odd $p^e\,\|\,N$ we have
$2^{s/D}\equiv1\pmod{p^e}$ once $\ord_{p^e}(2)\mid s/D$, so
$R_1(2^{s/D})\equiv R_1(1)$; for $2^e\,\|\,N$, $2^{s/D}\equiv0\pmod{2^e}$ once
$s/D\ge e$, so $R_1(2^{s/D})\equiv R_1(0)$.  This proves ``if''.  For the
converse and the refinement, let $s\to0$ in $\Zh$ along the set (e.g.\
$s=L\cdot\ord_{p^E}(2)\cdot D\cdot j$ with $E\to\infty$ and $j$ large); then
$2^{s/D}\to1$ in $\Z_p$ for odd $p$ and $2^{s/D}\to0$ in $\Z_2$, $R$ is
continuous on $\Q_p$, and $p^e\Z_p$ is closed.
\end{proof}

Thus a rule polynomial in $Y=2^{s/D}$ is valid near $\mu=0$ iff its
coordinates are positive for large $Y$ and their values at the two
specialisations $Y=1$ (the member $s=0$, $n=4^J$) and $Y=0$ (the
$2$-adic limit, $n=k$) are $\Z[1/2]$- resp.\ $\Z_{(2)}$-integral.  Neither
specialisation is locally obstructed; whether a rule exists is a global
question about polynomial identities.  For example, the Sophie Germain
rule of Theorem~\ref{thm:proth3} below specialises at $Y=X=1/2$ to
$x=K=3(4X^4+1)\mapsto 15/4$, $f_\pm=2X^2\pm2X+1\mapsto 5/2,\,1/2$: the
budget $K+d$ with $d=f_+^2f_-$ becomes $15/4+25/8=55/8$, divisible by
$11$ in $\Z[1/2]$.  The odd $2$-adic valuation of $f_\pm(1/2)$ is what makes
this possible (Lemma~\ref{lem:limitjacobi}).

\section{Theorems}\label{sec:theorems}

\subsection{The Proth family $3\cdot 2^s+1$}

\begin{theorem}[\status{proved}; {\cite[Thm.~C]{MenuSterility}}]\label{thm:proth3}
For every $s\ge0$, $4/(3\cdot2^s+1)$ is a sum of three unit fractions.
\end{theorem}

\begin{proof}
Write $n=3\cdot2^s+1$.  For $s=0$, $4/4=1/2+1/3+1/6$.

\emph{Step 1: the Sophie Germain window (for $4\mid s$, $s\ge4$).}  Put
$X=2^{(s-4)/4}$, so $n=48X^4+1$, and
\[
 K:=\frac{n+11}4=3(4X^4+1)=3f_+f_-,\qquad f_\pm:=2X^2\pm2X+1 .
\]
For $(\beta,\gamma)\in\{(2,1),(0,2),(2,0),(1,2)\}$ let $d=f_+^\beta f_-^\gamma$.
Then $d\mid K^2$ in $\Z[X]$, and $\gcd(11,K)=1$ because $4X^4\equiv-1\pmod{11}$
has no solution ($-1/4\equiv8$ is not a fourth power, indeed not a square,
modulo $11$).  By Lemma~\ref{lem:window}(b) with $x=K$, $m=11$,
\[
 x=K,\qquad y=\frac{n(K+d)}{11},\qquad z=\frac{nK(K+d)}{11\,d}
\]
is a solution as soon as $11\mid K+d$.  Since $K+d=f_+f_-(3+f_+)$,
$f_-(3f_++f_-)$, $f_+(3f_-+f_+)$, $f_+f_-(3+f_-)$ in the four cases and
$11\nmid f_+f_-$, the condition $11\mid K+d$ reads, respectively,
\begin{align*}
 (\beta,\gamma)=(2,1):&\ X^2+X+2\equiv0, & X&\equiv4,6 \pmod{11},\\
 (0,2):&\ 2X^2+X+1\equiv0, & X&\equiv2,3,\\
 (2,0):&\ 2X^2-X+1\equiv0, & X&\equiv8,9,\\
 (1,2):&\ X^2-X+2\equiv0, & X&\equiv5,7 .
\end{align*}
So every $X\not\equiv0,\pm1\pmod{11}$ is covered.  All of $n,K,f_\pm,d,K+d$
are positive for real $X$ (the discriminant of $f_\pm$ is $-4$), so the
solution is positive.  Since $2$ is a primitive root modulo $11$ and
$X=2^{(s-4)/4}$, $X\equiv1$ iff $s\equiv4\pmod{40}$ and $X\equiv-1$ iff
$s\equiv24\pmod{40}$.  Hence ES holds for all $s\equiv0\pmod4$ with
$s\not\equiv4,24\pmod{40}$.

\emph{Step 2: classes polynomial in $n$.}  Each line is an identity checked by
clearing denominators, valid with positive integer coordinates on the
stated class:
\begin{center}\small
\begin{tabular}{lll}
\toprule
condition on $n$ & exponents $s$ & $(x,y,z)$\\
\midrule
$7\mid n$ & $s\equiv1\pmod3$ & $(2n/7,\ 15n/7,\ 30n)$\\
$n\equiv6\pmod7$ & $s\equiv2\pmod3$ & $(2(n+1)/7,\ 2n,\ 2n(n+1)/7)$\\
$n\equiv7\pmod{15}$ & $s\equiv1\pmod4$ & $((4n+2)/15,\ 4n,\ 4n(2n+1)/15)$\\
$n\equiv13\pmod{15}$ & $s\equiv2\pmod4$ & $((4n+8)/15,\ 4n,\ 2n(n+2)/15)$\\
$5\mid n$ & $s\equiv3\pmod4$ & $(2n/5,\ 4n/5,\ 4n)$\\
$61\mid n$ & $s\equiv24\pmod{60}$ & $(16n/61,\ 326n/61,\ 2608n)$\\
\bottomrule
\end{tabular}
\end{center}
(The exponent classes follow from $\ord_7(2)=3$, $\ord_{15}(2)=4$,
$\ord_{61}(2)=60$ and $3\cdot2^{24}\equiv-1\pmod{61}$.)

\emph{Coverage.}  Let $s\ge1$.  If $3\nmid s$ or $4\nmid s$, a line of
Step~2 applies.  If $12\mid s$ and $s\not\equiv4,24\pmod{40}$, Step~1
applies.  If $12\mid s$ and $s\equiv4$ or $24\pmod{40}$, then
$s\equiv84$ or $24\pmod{120}$, so $s\equiv24\pmod{60}$ and $61\mid n$.
\end{proof}

\begin{remark}[structure and history]\label{rem:proth3}
(i) The Sophie Germain step is the only one that reaches the unit-square
point $\mu=0$: the class $s\equiv0\pmod{40}$ ($X\equiv6$) lies in Step~1.
By Theorem~\ref{thm:sterile} no finite set of classes polynomial in $n$ could
do this.  The two classes $X\equiv\pm1\pmod{11}$ that Step~1 misses are the
neighbourhoods of the integer point $X=\pm1$, where $48X^4+1=7^2$ is a
square; Step~2 handles them through $61$.  (A rule polynomial in $X$ with
integrality depending only on $X\bmod 33$ cannot be valid on a class
containing $X=\pm1$, by the odd-square argument applied to the integer
$49$, \cite[\S4.3]{MenuSterility}.)

(ii) A second, independent certificate was found in \cite[Thms.~2.2,
2.3]{Exp105}: for $3\cdot4^s+1$ a tree with nine leaves (three windows
$m=11$, two of Sophie Germain type), and for $6\cdot4^s+1$ the single class
I3$(1,5,3)$ for $s\ge1$.  It is re-verified by our script.  A third
mechanism reaches the same point: for $s\equiv0,12\pmod{36}$, $s\ge36$, with
$Y=2^{s/3}$, the Type~II family $a=80$, $c=1$, $d=(21Y+64)/320$,
$b=(Y^2-3Y+9)/7$ is valid.  Here $a,c$ are constant and the Type~II identity
$4abcd=cn+a+b$ is equivalent to $cn+a=b(4acd-1)$; the two factors are
$cn+a=3(Y+3)(Y^2-3Y+9)$ and $4acd-1=21(Y+3)$ \cite[Prop.~2.1]{ExpOpen}.

(iii) The family contains the primes $3\cdot2^s+1$ for
$s=1,2,5,6,8,12,18,30,36,41,66,189,\dots$.  For $s\ge3$ a member lies in one
of Mordell's six square classes modulo $840$ exactly when $12\mid s$; then
$n\equiv529\pmod{840}$ ($n\equiv1\pmod{24}$, $n\equiv4\pmod5$,
$n\equiv4\pmod7$).  So the content of Theorem~\ref{thm:proth3} beyond
Mordell's identities is the case $12\mid s$.
\end{remark}

\subsection{The family $63\cdot2^s+1$}

\begin{lemma}[\status{proved}; {\cite[Thm.~2.3]{ExpOpen}}]\label{lem:63class}
Let $s\equiv0\pmod{24}$, $s\ge24$, $Y=2^{s/3}$, $n=63Y^3+1=63\cdot2^s+1$,
$a=5$, $d=17$ and
\[
 b=\frac{63Y^3+3Y^2-9Y+28}{340},\qquad c=\frac{21Y+64}{340}.
\]
Then $(x,y,z)=(abd,\ acdn,\ bcdn)$ is a solution of \eqref{eq:ES} in positive
integers.  Its gap is $4x-n=3(Y^2-3Y+9)$.
\end{lemma}

\begin{proof}
We have $n+4a^2d=63Y^3+1701=63(Y+3)(Y^2-3Y+9)$.  Put $f=3(Y^2-3Y+9)$ and
$f'=21(Y+3)$, so $ff'=n+1700=n+4a^2d$.  One checks $340b-n=f$ and
$340c-1=f'$, i.e.\ $4abd-n=f$ and $4acd-1=f'$; so \eqref{eq:binom} holds
and $4abcd=cn+a+b$, which is the Type~II identity
(Lemma~\ref{lem:solve}).  Integrality: $340=4\cdot5\cdot17$; $4\mid Y$ since
$s\ge6$; $\ord_5(2)=4$ and $\ord_{17}(2)=8$, so $8\mid s/3$ gives
$Y\equiv1\pmod{85}$, and then $21Y+64\equiv85\equiv0$ and
$63Y^3+3Y^2-9Y+28\equiv85\equiv0\pmod{85}$, while both are $\equiv0\pmod4$.
Positivity is clear.
\end{proof}

The class $s\equiv0\pmod{24}$ contains a neighbourhood of the unit-square
point $\mu=0$ ($63+1=4^3$).  By Theorem~\ref{thm:sterile} no class polynomial
in $n$ reaches it, and by Theorem~\ref{thm:W} below no rule with constant
gap or constant ratio does either; the gap of Lemma~\ref{lem:63class} is the
cyclotomic factor $3(Y^2-3Y+9)$ of $Y^3+27$.  At $s=0$ the formulas give
$c=85/340$, so the condition $s\ge24$ is needed.

\begin{theorem}[\status{proved}, computer-assisted; {\cite[Thm.~3.1]{ExpOpen}}]\label{thm:63}
For every $s\ge0$, $4/(63\cdot2^s+1)$ is a sum of three unit fractions.
\end{theorem}

\begin{proof}
For $s\equiv0\pmod{24}$, $s\ge24$, use Lemma~\ref{lem:63class}.  The
eighteen ET and divisor classes of Table~\ref{tab:63} apply on the stated
exponent classes (for the classes of $n$ this follows from
\eqref{eq:period} and the listed orders; the identities are those of
Table~\ref{tab:ET}).  They cover every residue $s\bmod2520$ except
$s\equiv0,840,1680\pmod{2520}$, all of which are $\equiv0\pmod{24}$
(computer-checked coverage of $2520$ residues).  The thresholds in the table
are at most $5$, and every $s\ge1$ below the threshold of its class lies in
another class of the table without threshold (checked).  For $s=0$,
$4/64=1/32+1/64+1/64$.
\end{proof}
\begin{table}[ht]
\centering\small
\begin{tabular}{lll}
\toprule
class & condition on $n$ & exponents\\
\midrule
D$(5)$ & $5\mid n$ & $s\equiv 3\pmod{4}$ \\
II2$(1,2,15)$ & $n\equiv7\ (15)$ & $s\equiv 1\pmod{4}$ \\
II2$(2,2,15)$ & $n\equiv13\ (15)$ & $s\equiv 2\pmod{4}$ \\
II2$(5,3,119)$ & $n\equiv57\ (119)$ & $s\equiv 4\pmod{8}$ \\
II2$(5,2,39)$ & $n\equiv34\ (39)$ & $s\equiv 4\pmod{12}$ \\
II2$(1,5,19)$ & $n\equiv18\ (19)$ & $s\equiv 14\pmod{18}$ \\
I1$(1,146,39)$ & $n\equiv545\ (584)$ & $s\equiv 2\pmod{9}$, $s\ge5$ \\
D$(11)$ & $11\mid n$ & $s\equiv 2\pmod{10}$ \\
II2$(3,1,11)$ & $n\equiv8\ (11)$ & $s\equiv 4\pmod{10}$ \\
II2$(1,1,11)$ & $n\equiv7\ (11)$ & $s\equiv 6\pmod{10}$ \\
D$(29)$ & $29\mid n$ & $s\equiv 20\pmod{28}$ \\
II2$(1,11,87)$ & $n\equiv43\ (87)$ & $s\equiv 24\pmod{28}$ \\
D$(127)$ & $127\mid n$ & $s\equiv 1\pmod{7}$ \\
II2$(1,32,127)$ & $n\equiv126\ (127)$ & $s\equiv 2\pmod{7}$ \\
I1$(1,86,23)$ & $n\equiv321\ (344)$ & $s\equiv 4\pmod{14}$, $s\ge5$ \\
II1$(1,254,15)$ & $n\equiv1001\ (1016)$ & $s\equiv 1,5\pmod{7}$, $s\ge3$ \\
I1$(1,1405,511)$ & $n\equiv5109\ (5620)$ & $s\equiv 28\pmod{140}$, $s\ge4$ \\
I2$(2,43,151)$ & $n\equiv193\ (344)$, $n\equiv54\ (151)$ & $s\equiv 140\pmod{210}$, $s\ge3$ \\
\bottomrule
\end{tabular}
\caption{The ET/divisor classes for $63\cdot2^s+1$ (Theorem~\ref{thm:63}).  The
exponent classes are derived from the identities by the script (one full
period); ``$s\ge s_1$'' marks where the $2$-adic part of the class or
positivity needs $s$ large.}\label{tab:63}
\end{table}

\begin{remark}\label{rem:63}
Since $n_s\equiv1\pmod{168}$ and $n_s\equiv3\cdot2^s+1\pmod5$, $n_s$ lies in a
Mordell square class modulo $840$ (namely $169$) iff $4\mid s$; Mordell's
identities already handle $s\not\equiv0\pmod4$.  The essential content of
Theorem~\ref{thm:63} is the neighbourhood $s\equiv0\pmod{840}$ of $\mu=0$.
The only squares in the family are $n_0=8^2$ and $n_8=127^2$: if
$63\cdot2^s+1=v^2$ with $s\ge1$ then $v$ is odd, $s\ge2$, and with $v=2A+1$ we
get $A(A+1)=63\cdot2^{s-2}$; the odd one of the coprime factors $A,A+1$ is a
divisor $\delta$ of $63$ and the other is $(63/\delta)2^{s-2}$, and of the six
choices of $\delta$ only $\delta=63$, $2^{s-2}=64$ works.  The square $127^2$
lies in the divisor class $D_{127}$, as it must by Lemma~\ref{lem:T1}.
\end{remark}

\subsection{The family $7\cdot4^s+9$}

Here $7+9=4^2$, so $\mu=0$ is a unit-square point and
Theorem~\ref{thm:sterile} applies.  For even $s\ge2$ one has
$n\equiv121\pmod{840}$, so every member with even $s\ge2$ is Mordell-hard.
The windows used are $m=7\cdot4^jt-9$ with $(j,t)\in\{(1,1),(1,4),(2,1),(2,6)\}$,
i.e.\ $m\in\{19,103,663\}$.  For such a window and $s>j$,
\[
 x=\frac{n+m}4=7\cdot4^{j-1}\bigl(4^{s-j}+t\bigr),\qquad N=nx=4^{j-1}\cdot7\cdot n\cdot\bigl(4^{s-j}+t\bigr),
\]
and $4^{s-j}+t$ is split by (CY) or (SG): for $t=w^h$ and $h\mid s-j$, with
$Y=4^{(s-j)/h}$, into $Y+w$ ($h=1$), $(Y+1)(Y^2-Y+1)$ ($h=3$, $w=1$),
$(Y+1)(Y^2-Y+1)(Y^6-Y^3+1)$ ($h=9$, $w=1$); for $(j,t)=(1,1)$ and $s$ even,
with $Y=2^{(s-2)/2}$, into $f_+f_-$, $f_\pm=2Y^2\pm2Y+1$.  A rule is a choice
of $d_1,d_2$ among products of these pieces and of $2^i$ ($i\le 2j-2$), and it
applies by Lemma~\ref{lem:window}(a) whenever $m\mid d_1+d_2$.

\begin{theorem}[\status{proved}, computer-assisted; {\cite[Thm.~2.6]{Exp105}}]\label{thm:c7}
For every $s\ge0$, $4/(7\cdot4^s+9)$ is a sum of three unit fractions.
\end{theorem}

\begin{proof}
The $22$ window rules of Table~\ref{tab:c7w} and the $11$ ET/divisor
classes of Table~\ref{tab:c7et} apply on the listed exponent classes; their
union contains every residue modulo $9180=2^2\cdot3^3\cdot5\cdot17$
(computer-checked), and all thresholds are $\le3$.  For $s\in\{1,2\}$ the
rules apply directly, and for $s=0$, $4/16=1/8+1/12+1/24$.
\end{proof}

\begin{table}[ht]
\centering\small
\begin{tabular}{lllllll}
\toprule
$m$ & $j$ & type & $Y$ & $d_1$ & $d_2$ & exponents $s$\\
\midrule
$19$ & $1$ & CY & $4^{s-1}$ & $1$ & $(Y+1)$ & $6\pmod{9}$ \\
$19$ & $1$ & SG & $2^{(s-2)/2}$ & $f_-$ & $7$ & $18,26\pmod{36}$ \\
$19$ & $1$ & SG & $2^{(s-2)/2}$ & $f_+$ & $7$ & $0,8\pmod{36}$ \\
$19$ & $1$ & SG & $2^{(s-2)/2}$ & $1$ & $f_-$ & $8,32\pmod{36}$ \\
$103$ & $1$ & CY & $4^{s-1}$ & $(Y+4)$ & $7$ & $44\pmod{51}$ \\
$103$ & $1$ & CY & $4^{s-1}$ & $1$ & $n$ & $8\pmod{51}$ \\
$103$ & $1$ & CY & $4^{s-1}$ & $1$ & $(Y+4)$ & $41\pmod{51}$ \\
$103$ & $2$ & CY & $4^{s-2}$ & $1$ & $2\cdot 7\cdot n$ & $17\pmod{51}$ \\
$103$ & $2$ & CY & $4^{s-2}$ & $2^{2}$ & $7\cdot n$ & $47\pmod{51}$ \\
$103$ & $2$ & CY & $4^{s-2}$ & $2^{2}$ & $7\cdot (Y+1)$ & $38\pmod{51}$ \\
$103$ & $2$ & CY & $4^{s-2}$ & $2$ & $(Y+1)$ & $32\pmod{51}$ \\
$103$ & $2$ & CY$_{3}$ & $4^{(s-2)/3}$ & $1$ & $2\cdot 7\cdot n\cdot (Y^2-Y+1)$ & $128\pmod{153}$ \\
$103$ & $2$ & CY$_{3}$ & $4^{(s-2)/3}$ & $2^{2}\cdot (Y^2-Y+1)$ & $7\cdot n$ & $137\pmod{153}$ \\
$103$ & $2$ & CY$_{3}$ & $4^{(s-2)/3}$ & $2\cdot (Y^2-Y+1)$ & $7\cdot (Y+1)$ & $32,74\pmod{153}$ \\
$103$ & $2$ & CY$_{3}$ & $4^{(s-2)/3}$ & $2^{2}$ & $7\cdot (Y^2-Y+1)$ & $29\pmod{153}$ \\
$103$ & $2$ & CY$_{3}$ & $4^{(s-2)/3}$ & $2\cdot (Y+1)$ & $7\cdot (Y^2-Y+1)$ & $65,98\pmod{153}$ \\
$103$ & $2$ & CY$_{3}$ & $4^{(s-2)/3}$ & $n\cdot (Y+1)$ & $2^{2}\cdot 7$ & $56\pmod{153}$ \\
$103$ & $2$ & CY$_{3}$ & $4^{(s-2)/3}$ & $2$ & $n\cdot (Y^2-Y+1)$ & $11\pmod{153}$ \\
$103$ & $2$ & CY$_{3}$ & $4^{(s-2)/3}$ & $(Y^2-Y+1)$ & $n$ & $20\pmod{153}$ \\
$103$ & $2$ & CY$_{9}$ & $4^{(s-2)/9}$ & $1$ & $2\cdot 7\cdot n\cdot (Y+1)$ & $146,155\pmod{459}$ \\
$103$ & $2$ & CY$_{9}$ & $4^{(s-2)/9}$ & $1$ & $7\cdot n\cdot (Y^2-Y+1)$ & $128,308\pmod{459}$ \\
$663$ & $2$ & CY & $4^{s-2}$ & $2$ & $7\cdot (Y+6)$ & $10\pmod{12}$ \\
\bottomrule
\end{tabular}
\caption{Window rules for $7\cdot4^s+9$: $x=(n+m)/4$,
$y=nx(d_1+d_2)/(md_2)$, $z=nx(d_1+d_2)/(md_1)$ (Lemma~\ref{lem:window}(a));
$f_\pm=2Y^2\pm2Y+1$.  CY$_h$ means the cyclotomic splitting of
$4^{s-j}+1=Y^h+1$.  Each rule is valid for $s\ge2$ in its class (for
$j=1$ also $s=1$ where applicable).}\label{tab:c7w}
\end{table}

\begin{table}[ht]
\centering\small
\begin{tabular}{lll}
\toprule
class & condition on $n$ & exponents\\
\midrule
D$(11)$ & $11\mid n$ & $s\equiv 2\pmod{5}$ \\
D$(541)$ & $541\mid n$ & $s\equiv 110\pmod{270}$ \\
I3$(1,5,3)$ & $n\equiv17\ (20)$, $f\mid n^2+20$ & $s\equiv 1\pmod{2}$, $s\ge2$ \\
I3$(1,11,3)$ & $n\equiv41\ (44)$, $f\mid n^2+44$ & $s\equiv 2,4\pmod{5}$, $s\ge2$ \\
I3$(1,206,15)$ & $n\equiv809\ (824)$, $f\mid n^2+824$ & $s\equiv 50\pmod{102}$, $s\ge3$ \\
I3$(5,19,119)$ & $n\equiv261\ (380)$, $f\mid n^2+1900$ & $s\equiv 14,50,86,122,158\pmod{180}$, $s\ge2$ \\
II2$(1,2,31)$ & $n\equiv23\ (31)$ & $s\equiv 3\pmod{5}$ \\
II2$(1,8,1887)$ & $n\equiv1855\ (1887)$ & $s\equiv 30\pmod{36}$ \\
II2$(2,2,111)$ & $n\equiv79\ (111)$ & $s\equiv 12\pmod{18}$ \\
II2$(7,2,671)$ & $n\equiv279\ (671)$ & $s\equiv 26\pmod{30}$ \\
II2$(19,6,455)$ & $n\equiv436\ (455)$ & $s\equiv 4\pmod{6}$ \\
\bottomrule
\end{tabular}
\caption{ET/divisor classes for $7\cdot4^s+9$; for I3 the modulus $f$ is the
third parameter.}\label{tab:c7et}
\end{table}

\begin{remark}\label{rem:c7}
The odd exponents are all handled by I3$(1,5,3)$; the work is in the even
exponents.  The rule that reaches $\mu=0$ is the Sophie Germain window
$m=19$ with $(d_1,d_2)=(f_+,7)$ on $s\equiv0,8\pmod{36}$: at the
specialisation $Y=1/2$ one has $f_+=5/2$ and $7+5/2=19/2$.  By
Lemma~\ref{lem:limitjacobi} below, only a piece with odd $2$-adic valuation
at the limit can do this, and $f_\pm(1/2)\in\{5/2,1/2\}$ are such pieces.
The tree that produced the certificate splits first on $s\bmod 36$ (the
period of the $m=19$ rule) and then on the primes $17$, $5$, $3$; the
certificate was replayed by two independent checkers in \cite{Exp105}.
\end{remark}

\subsection{The family $4^s+105$}
Since $105=3\cdot5\cdot7\equiv1\pmod 8$, $n_s=4^s+105\equiv(2^s)^2\pmod{105}$ and
$n_s\equiv1\pmod8$ for $s\ge2$: every member with $s\ge2$ lies in one of
Mordell's six square classes modulo $840$, and $n_s$ is a square modulo every
prime factor of the offset.  The family has the non-unit square points
$4^2+105=11^2$, $4^3+105=13^2$, $4^4+105=19^2$ and $1+105\cdot4^2=41^2$
($\mu=-2$), and a greedy search for an ET covering stalled near ``mixed''
points of $\Zh$ that imitate different integer squares at different
primes \cite[\S1.4]{SignFields}.  One cubic window class closes it.

\begin{lemma}[\status{proved}; {\cite[Lemmas~1.0, 1.0$'$]{Sparks6}}]\label{lem:105class}
Let $u\equiv32\pmod{151}$ be a positive integer, $n=256u^3+105$ and
\[
 x=64(u^3+1),\qquad A=\frac{64u^3+2u^2-2u+66}{151},\qquad y=nA,\qquad z=32n(u+1)A .
\]
Then $A\in\N$ and $4/n=1/x+1/y+1/z$.  If $s\equiv34\pmod{45}$ then
$u:=4^{(s-4)/3}$ is such an integer and $n=4^s+105$.
\end{lemma}

\begin{proof}
$4x-n=151$, so $4/n-1/x=151/(nx)$, and
$1/y+1/z=(32u+33)/(32n(u+1)A)$; equality with $151/(64n(u^3+1))$ is
$151A=2(32u+33)(u^2-u+1)$, which is the definition of $A$ expanded.  For
$u\equiv32$, $32u+33\equiv1057=7\cdot151\equiv0\pmod{151}$, so
$A\in\N$.  This is Lemma~\ref{lem:window}(a) with $m=151=4^4-105$,
$N=nx$, $d_1=32(u+1)$ and $d_2=1$\,---\,more transparently,
Lemma~\ref{lem:window}(b) with $d=x/(32(u+1))=2(u^2-u+1)$, a divisor of
$x$ forced by $u^3+1=(u+1)(u^2-u+1)$.  For the second statement,
$\ord_{151}(4)=15$, $(s-4)/3\equiv10\pmod{15}$ and $4^{10}\equiv32\pmod{151}$.
\end{proof}

\begin{theorem}[\status{proved}, computer-assisted; {\cite[Thm.~1.1]{Sparks6}}]\label{thm:105}
For every $s\ge0$, $4/(4^s+105)$ is a sum of three unit fractions.
\end{theorem}

\begin{proof}
The thirteen ET/divisor classes of Table~\ref{tab:105} cover every residue
$s\bmod180$ except $s\equiv34\pmod{90}$ (computer-checked, or by hand from the
table: the classes modulo $5$, $10$, $20$, $6$, $9$, $15$, $30$ leave
$s\equiv34,124\pmod{180}$), and $s\equiv34\pmod{90}$ implies
$s\equiv34\pmod{45}$, where Lemma~\ref{lem:105class} applies.  All moduli are
odd and all coefficients are non-negative, so the classes apply for every
$s\ge0$.
\end{proof}

\begin{table}[ht]
\centering\small
\begin{tabular}{lll}
\toprule
class & condition on $n$ & exponents\\
\midrule
D$(11)$ & $11\mid n$ & $s\equiv 2\pmod{5}$ \\
II2$(1,1,11)$ & $n\equiv7\ (11)$ & $s\equiv 0\pmod{5}$ \\
II2$(1,3,11)$ & $n\equiv10\ (11)$ & $s\equiv 1\pmod{5}$ \\
D$(41)$ & $41\mid n$ & $s\equiv 8\pmod{10}$ \\
II2$(2,53,1271)$ & $n\equiv423\ (1271)$ & $s\equiv 9\pmod{10}$ \\
II2$(3,22,527)$ & $n\equiv262\ (527)$ & $s\equiv 13\pmod{20}$ \\
II2$(1,6,143)$ & $n\equiv119\ (143)$ & $s\equiv 24\pmod{30}$ \\
II2$(1,2,151)$ & $n\equiv143\ (151)$ & $s\equiv 14\pmod{15}$ \\
D$(13)$ & $13\mid n$ & $s\equiv 3\pmod{6}$ \\
D$(19)$ & $19\mid n$ & $s\equiv 4\pmod{9}$ \\
II2$(5,1,19)$ & $n\equiv14\ (19)$ & $s\equiv 1\pmod{9}$ \\
II2$(2,5,39)$ & $n\equiv37\ (39)$ & $s\equiv 5\pmod{6}$ \\
II2$(1,2,39)$ & $n\equiv31\ (39)$ & $s\equiv 1\pmod{6}$ \\
\bottomrule
\end{tabular}
\caption{ET/divisor classes for $4^s+105$ (Theorem~\ref{thm:105}).}\label{tab:105}
\end{table}

\begin{remark}\label{rem:105}
(i) The window class respects the odd-square obstruction locally: for
$q=18837001$, a prime with $\ord_q(4)=45$, one has $\leg{n_s}q=-1$ for every
$s\equiv34\pmod{45}$ \cite[review R246]{Sparks6}.  (ii) Whether
ET and divisor classes alone cover $4^s+105$ is \status{open}: no point of
the closure of the family is known to be sterile, the greedy search merely
stalled (uncovered Haar measure $8.5\cdot10^{-8}$ of the exponent space,
\status{evidence}).  So the theorem bypasses, and does not refute, a possible
ET covering.  (iii) $4^s+105$ is not the first all-Mordell-hard family
settled: $315\cdot2^m+1\equiv1\pmod{2520}$ for $m\ge3$ was settled by an ET
covering in \cite{ExpCov}.  The new feature is a class whose divisor
$2(u^2-u+1)$ moves with $n$.  The architecture\,---\,a finite exponent
covering leaving one class, closed by an algebraic factorisation along that
class\,---\,is Izotov's \cite{Izotov,FFK}.
\end{remark}

\subsection{The families $4^s+k$}
\begin{theorem}[\status{proved}, computer-assisted; {\cite[Thm.~1.3]{SignFields}}]\label{thm:nine}
For $k\in\{9,57,81,129,153,177,201,225,249\}$ and every $s\ge0$,
$4/(4^s+k)$ is a sum of three unit fractions.
\end{theorem}

\begin{proof}
For each $k$ the ET/divisor classes of Table~\ref{tab:nine} apply on the
listed exponent classes for all $s\ge0$ (for I1$(1,679,19)$, $s\ge2$; the
member $s=1$, $n=229$, of that class is covered by II2$(1,1,7)$), and their
union is all residues modulo the stated period.
\end{proof}

\begin{table}[ht]
\centering\small
\begin{tabular}{llll}
\toprule
$k$ & class & condition on $n$ & exponents\\
\midrule
$9$ & D$(5)$ & $5\mid n$ & $s\equiv 0\pmod{2}$ \\
 & II2$(2,2,15)$ & $n\equiv13\ (15)$ & $s\equiv 1\pmod{2}$ \\
\midrule
$57$ & II2$(2,2,15)$ & $n\equiv13\ (15)$ & $s\equiv 0\pmod{2}$ \\
 & II2$(2,1,7)$ & $n\equiv5\ (7)$ & $s\equiv 1\pmod{3}$ \\
 & II2$(1,1,7)$ & $n\equiv3\ (7)$ & $s\equiv 2\pmod{3}$ \\
 & D$(11)$ & $11\mid n$ & $s\equiv 3\pmod{5}$ \\
 & II2$(1,1,11)$ & $n\equiv7\ (11)$ & $s\equiv 2\pmod{5}$ \\
 & II2$(1,1,31)$ & $n\equiv27\ (31)$ & $s\equiv 0\pmod{5}$ \\
 & II2$(1,8,31)$ & $n\equiv30\ (31)$ & $s\equiv 1\pmod{5}$ \\
 & II2$(7,9,2015)$ & $n\equiv251\ (2015)$ & $s\equiv 9\pmod{30}$ \\
\midrule
$81$ & D$(5)$ & $5\mid n$ & $s\equiv 1\pmod{2}$ \\
 & II2$(1,2,15)$ & $n\equiv7\ (15)$ & $s\equiv 0\pmod{2}$ \\
\midrule
$129$ & D$(5)$ & $5\mid n$ & $s\equiv 0\pmod{2}$ \\
 & II2$(2,2,15)$ & $n\equiv13\ (15)$ & $s\equiv 1\pmod{2}$ \\
\midrule
$153$ & II2$(1,2,15)$ & $n\equiv7\ (15)$ & $s\equiv 1\pmod{2}$ \\
 & D$(7)$ & $7\mid n$ & $s\equiv 0\pmod{3}$ \\
 & II2$(1,1,7)$ & $n\equiv3\ (7)$ & $s\equiv 1\pmod{3}$ \\
 & D$(13)$ & $13\mid n$ & $s\equiv 2\pmod{6}$ \\
\midrule
$177$ & II2$(2,2,15)$ & $n\equiv13\ (15)$ & $s\equiv 0\pmod{2}$ \\
 & II2$(1,1,7)$ & $n\equiv3\ (7)$ & $s\equiv 0\pmod{3}$ \\
 & II2$(1,2,7)$ & $n\equiv6\ (7)$ & $s\equiv 1\pmod{3}$ \\
 & II2$(1,2,39)$ & $n\equiv31\ (39)$ & $s\equiv 5\pmod{6}$ \\
\midrule
$201$ & D$(5)$ & $5\mid n$ & $s\equiv 1\pmod{2}$ \\
 & II2$(1,2,15)$ & $n\equiv7\ (15)$ & $s\equiv 0\pmod{2}$ \\
\midrule
$225$ & II2$(2,1,7)$ & $n\equiv5\ (7)$ & $s\equiv 1\pmod{3}$ \\
 & II2$(1,1,7)$ & $n\equiv3\ (7)$ & $s\equiv 2\pmod{3}$ \\
 & II2$(1,2,39)$ & $n\equiv31\ (39)$ & $s\equiv 0\pmod{6}$ \\
 & D$(17)$ & $17\mid n$ & $s\equiv 3\pmod{4}$ \\
 & II2$(101,5,2019)$ & $n\equiv1918\ (2019)$ & $s\equiv 9\pmod{24}$ \\
 & I1$(1,679,19)$ & $n\equiv2697\ (2716)$ & $s\equiv 21\pmod{24}$, $s\ge2$ \\
\midrule
$249$ & D$(5)$ & $5\mid n$ & $s\equiv 0\pmod{2}$ \\
 & II2$(2,2,15)$ & $n\equiv13\ (15)$ & $s\equiv 1\pmod{2}$ \\
\bottomrule
\end{tabular}
\caption{ET/divisor classes for $4^s+k$ (Theorem~\ref{thm:nine}).}\label{tab:nine}
\end{table}

\begin{remark}[scope]\label{rem:nine}
By the elementary reduction \cite[Lemma~1.1]{SignFields}, $4^s+k$ with
$s\ge2$ can meet Mordell's square classes only if $k\equiv9\pmod{24}$, which
holds for all nine $k$ and for $33,105,273,297$ among $k\le300$.  For
$k\in\{9,81,129,201,249\}$ no member with $s\ge2$ lies in a Mordell square
class, so these cases are Mordell's theorem; $k=153$ is Mordell's theorem
plus the divisor class $13\mid n$.  The genuinely new coverings are
$k=57$, $177$, $225$.  The families $4^s+k$ have no integer unit-square
points (Proposition~\ref{prop:intpoints}), so Theorem~\ref{thm:sterile} does
not obstruct them; the covering is of Sierpi\'nski type.  For
$k\in\{33,273,297\}$ ES for all $4^s+k$ is \status{open}: tree searches with
ET classes, windows and sums-of-powers windows stall at uncovered measures
$\approx10^{-7}$ near the square points $4^2+33=7^2$, $4^4+273=23^2$,
$4^3+297=19^2$ \cite{SignFields,Sparks6,Exp105,ExpOpen}.
\end{remark}
